\documentclass[10pt,aspectratio=169]{beamer}

% ---------------------------------------------------------------
% Theme and typography (XeLaTeX + metropolis defaults: Fira Sans)
% ---------------------------------------------------------------
\usetheme{metropolis}
\metroset{progressbar=frametitle,numbering=fraction,block=fill,sectionpage=progressbar}

% Suppress the bookmarks/outline pane on PDF open (user preference;
% see ~/.claude/projects/C--git-ckt/memory/feedback_pdf_no_bookmarks.md)
\hypersetup{pdfpagemode=UseNone}

\usepackage{fontspec}
\usepackage{amsmath,amsthm,mathtools,bbm}
% unicode-math replaces amssymb and matches math letters to a font
% that visually pairs with Metropolis's Fira Sans (math letters in
% the default Computer Modern look heavy and oversized next to Fira
% Sans body text).
\usepackage{unicode-math}
% Fira Math would visually pair with Fira Sans (Metropolis's text font),
% but it's not bundled with MiKTeX. Fall back to Latin Modern Math:
% serif, but proper Unicode math with consistent sizing across script
% and scriptscript styles --- much better than the Computer Modern
% default which produces oversized math letters.
\setmathfont{Latin Modern Math}
\usepackage{booktabs}
\usepackage{tikz}
\usetikzlibrary{arrows.meta,positioning,calc,decorations.pathreplacing}
\usepackage{xcolor}
\usepackage[authoryear,round]{natbib}
% Suppress the box that beamer puts around bib items by default.
\setbeamertemplate{bibliography item}{}

% CKT-style colors
\definecolor{cktblue}{HTML}{1F4E79}
\definecolor{cktaccent}{HTML}{C0504D}
\definecolor{cktgray}{HTML}{595959}
\definecolor{cktgreen}{HTML}{2E7D32}
\setbeamercolor{frametitle}{bg=cktblue,fg=white}
\setbeamercolor{progress bar}{fg=cktaccent,bg=cktblue!20}
\setbeamercolor{title separator}{fg=cktaccent}
\setbeamercolor{alerted text}{fg=cktaccent}

% Notation shortcuts matching CKT
\newcommand{\D}{D}
\newcommand{\thetai}{\theta_{i}}
\newcommand{\Deltai}{\Delta_{i}}
\newcommand{\phicoef}{\phi}
\newcommand{\mud}{\mu_{\underline{d}}}
\newcommand{\Deltad}{\Delta_{\underline{d}}}
\newcommand{\DS}{\mathcal{D}_{S}}
\newcommand{\dN}{d_{N}}
\newcommand{\dT}{d_{T}}
\newcommand{\vii}{v_{i}}
\newcommand{\nuil}{\nu_{it}^{l}}
\newcommand{\tildenu}{\tilde{\nu}_{it}^{l}}

\title{Implementing Verdier's robust extrapolation in the CKT model}
\subtitle{Relaxing the i.i.d.\ assumption on $\nuil$}
\author{Cenci, Kleemans, Tjernstr\"om}
\date{\today}

\begin{document}

% ---------------------------------------------------------------
\begin{frame}[plain]
  \titlepage
\end{frame}

% ---------------------------------------------------------------
% Slide 2
\begin{frame}{Our current estimator corresponds to Verdier's simple extrapolation}

  We currently follow \citet{suriSelectionComparativeAdvantage2011} and \citet{tjernstromCommentSuri2011}:
  \begin{itemize}
    \item impose the LCA restriction $\Deltai = \beta + \phicoef\, \thetai$
    \item identify $\phicoef$ from switcher trajectories
    \item extrapolate to non-migrants ($\dN$, $\dT$) along the LCA line
  \end{itemize}

  \vspace{0.8em}
  \citet{verdierAverageTreatmentEffects2020a} treats this class of methods as a ``simple extrapolation'' \& proposes a robust alternative.
  Our estimator and his simple version are not literally the same procedure (different aggregation level, reciprocal slope convention) but target the same LCA relationship asymptotically.

\end{frame}

% ---------------------------------------------------------------
% Slide 3
\begin{frame}{The simple extrapolation assumes i.i.d.\ non-pecuniary shocks}

  Our current Assumption A2 states that the non-pecuniary utility shock
  \[
    \nuil \quad \text{is i.i.d.\ across individuals } i, \text{ time } t, \text{ \& locations } l
  \]

  \vspace{0.6em}
  This implies:
  \begin{itemize}
    \item $\nuil$ independent of baseline skill \& comparative advantage \alert{everywhere}
    \item rules out systematic differences in $\nuil$ across regions/institutions
    \item workers in different locations have the same shock distribution, conditional on comparative advantage
  \end{itemize}

\end{frame}

% ---------------------------------------------------------------
% Slide 4: finding 15 --- drop the China bullet
\begin{frame}{The migration literature documents what A2 rules out}

  Regional heterogeneity is a central fact of this literature:

  \vspace{0.6em}
  \begin{itemize}
    \item \citet{bryanAggregateProductivityEffects2019}: structural model of Indonesian internal migration estimates large bilateral migration costs
    \item \citet{lagakosMigrationCostsObservational2020}: migration cost estimates differ across developing countries
    \item \citet{pulidoBarriersMobilitySorting2021}: cross-sector barriers in Indonesia explain a large share of within-country income gaps
  \end{itemize}

  \vspace{0.6em}
  Our current model rules this out by assumption

\end{frame}

% ---------------------------------------------------------------
% Slide 5: findings 5, 6 --- replace alertblock; drop bold pseudo-label
\begin{frame}{Verdier's robust extrapolation weakens A2 to a conditional version}

  Decompose the non-pecuniary shock:
  \[
    \nuil \;=\; m_{l}(\vii) \;+\; \tildenu
  \]
  where $\vii$ is an observed, time-invariant individual-level index

  \vspace{0.6em}
  New assumption (A2$'$):
  \begin{itemize}
    \item $\tildenu$ is i.i.d.\ \emph{conditional on} $\vii$
    \item $m_{l}(\vii)$ can correlate arbitrarily with comparative advantage
  \end{itemize}

  \vspace{0.6em}
  We use first-wave province as $\vii$ to index region-level shocks (all three countries)

  \vspace{0.6em}
  \textit{A2 (our assumption) requires unconditional independence} \\
  \textit{A2$'$ (VV's assumption) requires only within-province independence}

\end{frame}

% ---------------------------------------------------------------
% Slide 6 (keep \textbf{Simple:}/\textbf{Robust:} --- they earn their weight)
% finding 9: strip accumulated \vspace calls
\begin{frame}{In practice: the LCA gets a cluster-specific intercept}

  Slope stays common but the intercept varies

  \vspace{0.4em}
  \textbf{Simple:}
  \[
    \Deltai \;=\; \beta \;+\; \phicoef\, \thetai
  \]

  \textbf{Robust:}
  \[
    \Deltai \;=\; \beta(\vii) \;+\; \phicoef\, \thetai
  \]

  \begin{itemize}
    \item slope $\phicoef$ remains common; this is still the LCA restriction
    \item intercept $\beta(\vii)$ absorbs cluster-level $m_l(\vii)$ heterogeneity
    \item $\phicoef$ identifies off within-cluster switcher variation
    \item extrapolation to non-migrants happens cluster-by-cluster, then averaged
  \end{itemize}

  $\vii$ is generic in VV's framework (village in his Suri-Kenya application; for CKT, see slide on geography below).

\end{frame}

% ---------------------------------------------------------------
% Slide 7 (keep \textbf{Province A}/\textbf{Province B} --- definitional contrast)
% finding 9: strip extra \vspace
\begin{frame}{Pooling provinces mixes differentially-selected switcher samples}

  Consider two provinces:

  \begin{itemize}
    \item \textbf{Province A} (many movers): wide switcher sample, broad range of $\thetai$
    \item \textbf{Province B} (few movers): skewed sample, only high-return types migrate
  \end{itemize}

  \vspace{0.6em}
  Pooled regression: can be biased by across-province difference in selection

  Within-province: clean slope, holds region-level shocks constant

\end{frame}

% ---------------------------------------------------------------
% Slide 8: findings 12, 13
% nudge green label up; convert flushleft block to itemize
\begin{frame}{Within-province slopes agree---the pooled slope is pulled off}

  \centering
  \begin{tikzpicture}[x=0.85cm,y=0.75cm,>={Stealth[length=2mm]}]

    % axes
    \draw[->,thick] (-0.2,0) -- (6.5,0) node[right] {\small $\thetai$};
    \draw[->,thick] (0,-0.2) -- (0,5.7) node[above] {\small $\Deltai$};

    % Province A: many movers, wide theta range, slope = phi (negative)
    % Points at slope -0.5: (1.0,3.0)+(0.8,-0.4) per step
    \foreach \x/\y in {1.0/3.0, 1.8/2.6, 2.6/2.2, 3.4/1.8, 4.2/1.4, 5.0/1.0} {
      \filldraw[cktblue,opacity=0.85] (\x,\y) circle (2pt);
    }
    \draw[cktblue,thick] (0.5,3.25) -- (5.5,0.75);

    % Province B: few movers concentrated at LOW theta, HIGH Delta
    % (when phi<0, high-return types are low-theta, so the skewed cloud sits upper-left)
    % Points at slope -0.5: (0.5,5.0)+(0.5,-0.25) per step
    \foreach \x/\y in {0.5/5.0, 1.0/4.75, 1.5/4.5, 2.0/4.25, 2.5/4.0} {
      \filldraw[cktgreen,opacity=0.85] (\x,\y) circle (2pt);
    }
    \draw[cktgreen,thick] (0.2,5.15) -- (3.0,3.75);

    % Pooled regression line: biased steeper than phi by between-province intercept gap.
    % Actual OLS slope across the 11 points is -0.861; intercept 5.132.
    % Drawn from (0.2, 4.96) to (5.5, 0.40) to match the OLS through the joint centroid.
    \draw[cktaccent,thick,dashed] (0.2,4.96) -- (5.5,0.40);

    % Legend (kept inside axis box at x ~ 6.0 to avoid clipping at x=6.5)
    \node[anchor=west, cktblue]   at (3.8,5.4) {\small Province $A$ (many movers)};
    \node[anchor=west, cktgreen]  at (3.8,4.9) {\small Province $B$ (few movers)};
    \node[anchor=west, cktaccent] at (3.8,4.4) {\small Pooled fit};

    % Slope annotation (one label, in clear space between the two parallel solid fit lines)
    \node[cktblue,font=\footnotesize]    at (3.4,3.1) {slope $= \phicoef$};
    \node[cktaccent,font=\footnotesize]  at (4.7,2.5) {biased fit};

  \end{tikzpicture}

  \vspace{0.4em}
  \begin{itemize}
    \item Each province's switcher cloud has true slope $\phicoef$
    \item The pooled line tilts because of intercept differences, not because $\phicoef$ varies
  \end{itemize}

\end{frame}

% ---------------------------------------------------------------
\section{How Suri, Verdier, and CKT identify the LCA slope}

% ---------------------------------------------------------------
% Slide 9 (keep \textbf{Stage 1}/\textbf{Stage 2} --- definitional)
% finding 9: strip \vspace under 0.6em where possible
\begin{frame}{Suri (2011) identifies $\phicoef$ via panel SUR \& minimum distance}

  \citet{suriSelectionComparativeAdvantage2011}'s two-stage indirect approach on a multi-year, binary-adoption panel:

  \textbf{Stage 1 (reduced form).}
  Stack one outcome equation per year. SUR over the panel:
  \[
    y_{it} \;=\; \pi_{0,t} \;+\; \sum_{\underline{d}\in\mathcal{D}} \pi_{\underline{d},t}\, \mathbbm{1}\{\underline{d}_i = \underline{d}\} \;+\; x_{it}'\gamma_t \;+\; u_{it}
  \]
  Each coefficient $\pi_{\underline{d}, t}$ is the average $y$ in adoption pattern $\underline{d}$ in year $t$

  \textbf{Stage 2 (structural / OMD).}
  Express each $\pi_{\underline{d}, t}$ as a known function of $(\beta, \phicoef, \lambda_1, \ldots, \lambda_K)$ implied by the CRC restriction
  \[
    g_{\underline{d}, t}(\beta, \phicoef, \lambda) \;=\; \pi_{\underline{d}, t}
  \]
  Solve via weighted minimum distance: $\hat\theta = \arg\min (\hat\pi - g(\theta))' W (\hat\pi - g(\theta))$

  \vspace{0.2em}
  \begin{itemize}
    \item Identification: SUR over-identifies the structural parameters; $\phicoef$ gets pinned by the across-equation cross-restrictions
    \item Implemented in our \texttt{randcoef.ado}, which replicates Suri's estimation approach
  \end{itemize}

\end{frame}

% ---------------------------------------------------------------
% Slide 10: finding 10 --- re-title as assertion; finding 2 --- remove \footnotesize
\begin{frame}{Step 1 estimates worker-level returns by saturated FE}

  \citet{verdierAverageTreatmentEffects2020a}'s CRC model (his eq.\ 2.4):
  \[
    y_{it} \;=\; a_i \;+\; b_i\, \D_{it} \;+\; f_t \;+\; x_{it}'\gamma \;+\; u_{it}
  \]
  Worker-level intercept $a_i$ (rural mean); worker-level treatment effect $b_i$ (return).

  \vspace{0.3em}
  Step 1 estimates this by OLS on $n$ worker indicators (absorbing $a_i$) plus $n$ worker-by-treatment interactions (recovering $b_i$), with period indicators / covariates.

  \vspace{0.3em}
  For each mover, this delivers $\hat a_i$ and $\hat a_i + \hat b_i$ as the within-$i$ regime means of $(y - x'\hat\gamma)$:
  \begin{align*}
    \hat a_i &= \text{avg of }(y - x'\hat\gamma)\text{ over }\D_{it} = 0\text{ periods}, \\
    \hat a_i + \hat b_i &= \text{avg of }(y - x'\hat\gamma)\text{ over }\D_{it} = 1\text{ periods}.
  \end{align*}
  Then $\hat b_i = (\hat a_i + \hat b_i) - \hat a_i$.

  \vspace{0.2em}
  Movers have both; untreated stayers only $\hat a_i$; treated stayers only $\hat a_i + \hat b_i$.

\end{frame}

% ---------------------------------------------------------------
% Slide 11: finding 10 --- re-title as assertion; finding 2 --- remove \footnotesize
\begin{frame}{Step 2 IV-regresses worker rural means on worker returns}

  VV writes the LCA in the direction $a$-on-$b$:
  \[
    a_i \;=\; \alpha_0 \;+\; \alpha_1\, b_i \;+\; e_i
  \]
  Naive OLS is biased: $\hat b_i$ is noisy, so it is a regressor with measurement error.

  \vspace{0.4em}
  VV instruments $\hat b_i$ with the binary trajectory pattern $(x_{i1}, \ldots, x_{iT})$:
  \[
    E\!\left[\, x_{it}\;\bigl(\hat a_i - \alpha_0 - \alpha_1 \hat b_i\bigr)\,\right] \;=\; 0, \quad t = 1, \ldots, T
  \]

  \vspace{0.4em}
  Robust version (cluster-level cost shifters): add cluster-$\vii$ FE to the IV regression.
  By Frisch-Waugh, demean LHS, RHS, \emph{and} instrument within $\vii$; intercept drops out:
  \[
    E\!\left[\, (x_{it} - \bar x_{\vii,t})\;\bigl((\hat a_i - \bar{\hat a}_{\vii}) - \alpha_1 (\hat b_i - \bar{\hat b}_{\vii})\bigr)\,\right] \;=\; 0
  \]

  \vspace{0.3em}
  Run on movers only.
  Cluster-robust SE at $\vii$.

\end{frame}

% ---------------------------------------------------------------
% Slide 12: finding 6 --- drop \textbf{Robust variant:}
\begin{frame}{CKT identifies $\phicoef$ at the trajectory level}

  Trajectory pooling: every worker with $\underline{d}_i = \underline{d}$ has the same expected baseline
  \[
    E[\thetai \mid \underline{d}_i = \underline{d}] \;=\; \mud - \mu_{\underline{d}_0}
  \]

  \vspace{0.3em}
  Single GMM equation, one scalar $\mud$ per trajectory:
  \begin{align*}
    y_{it} \;=\; \sum_{\underline{d}} \mud \mathbbm{1}\{\underline{d}_i = \underline{d}\}
      \;+\; \Delta_{\underline{d}_0} \D_{it}
      \;+\; \phicoef \!\sum_{\underline{d}\in\DS\setminus\underline{d}_0}\!\! (\mud - \mu_{\underline{d}_0}) \D_{it} \mathbbm{1}\{\underline{d}_i = \underline{d}\}
      \;+\; \cdots
  \end{align*}

  \vspace{0.3em}
  Identification: switcher$\times$choice indicators as instruments

  \vspace{0.6em}
  The robust variant uses the same equation with switcher$\times$choice instruments village-demeaned (like VV)

\end{frame}

% ---------------------------------------------------------------
% Slide 13: finding 7 --- remove trailing periods from two closing bullets
\begin{frame}{Both procedures collapse to the same trajectory-level moments}

  VV's IV groups movers by trajectory pattern $g$ and uses within-group means $(\bar{\hat a}_g, \bar{\hat b}_g)$.
  Within-group averaging washes out the Step 1 measurement error in $\hat b_i$.
  With $T = 2$ and two mover bins $(0,1)$, $(1,0)$, the IV reduces to a Wald estimator:
  \[
    \hat\alpha_1 \;=\; \frac{\bar{\hat a}_{(0,1)} - \bar{\hat a}_{(1,0)}}{\bar{\hat b}_{(0,1)} - \bar{\hat b}_{(1,0)}}
  \]

  \vspace{0.3em}
  CKT computes per-trajectory $\mud$ and $\Deltad$ directly via GMM.
  Same set of trajectory-level points; both procedures fit a line through them.

  \vspace{0.4em}
  \begin{itemize}
    \item Without cluster structure: VV-simple and CKT-simple are equivalent up to slope-direction convention
    \item With cluster structure: CKT-robust pools across clusters within trajectory (assumption A3 below)
    \item VV-robust uses only within-cluster variation; no across-cluster pooling needed
  \end{itemize}

\end{frame}

% ---------------------------------------------------------------
% Breath slide between slides 13 and 14
\begin{frame}[c]{}
  \centering\large
  How do $\alpha_1$ and $\phicoef$ relate?
\end{frame}

% ---------------------------------------------------------------
% Slide 14: finding 6 --- drop \textbf{Regression direction.} and
% \textbf{Right-hand-side variable.}; rephrase as First/Second prose;
% finding 2 --- remove \footnotesize
\begin{frame}{VV's $\alpha_1$ and CKT's $\phicoef$ target different population scalars}

  VV's LCA: $a_i = \alpha_0 + \alpha_1 b_i + e_i$ \qquad ($\alpha_1$: slope of $a$ on $b$)

  CKT's LCA: $\Deltai = \beta + \phicoef \thetai + \xi_i$ \qquad ($\phicoef$: slope of $\Delta$ on $\theta$)

  \vspace{0.4em}
  Two distinctions:
  \begin{itemize}
    \item First, the regression directions differ: slope of $a$ on $b$ vs slope of $\Delta$ on $\theta$.
          Reciprocal only when $|\mathrm{Corr}(a_i, b_i)| = 1$ (perfect linear LCA).
          Sign agrees automatically when both target the same correlation; magnitudes differ.
    \item Second, the right-hand-side variables differ: rural baseline $a_i$ vs comparative advantage $\thetai = b_R(\theta_i^U - \theta_i^R)$.
          CKT's structural model interprets $\Deltai$ in terms of two skills $(\theta_i^U, \theta_i^R)$; VV is agnostic about that structure.
          Both target the same scalar empirically (under correct assumptions); CKT's two-skill structure buys interpretation, not identification.
  \end{itemize}

  \vspace{0.4em}
  See appendix for the algebra.

\end{frame}

% ---------------------------------------------------------------
% Finding 1 & 2: slide 15 was "Three estimators" with \footnotesize covering
% both assumptions and the table. Split into two slides.
%
% Slide 15a: Assumptions A1--A3 (description-style itemize, no \footnotesize)
% Keep \textbf{A1}/\textbf{A2}/\textbf{A2'}/\textbf{A3} labels --- they earn weight
% Finding 7: remove trailing periods from all four description items
\begin{frame}{A2$'$ is strictly weaker than A2; A3 is the cost of trajectory pooling}

  \begin{description}
    \item[\textbf{A1}] LCA at the worker level: $\Deltai = \beta(\vii) + \phicoef\,\thetai + \xi_i$
    \item[\textbf{A2}] Unconditional: $\nuil$ i.i.d.\ across workers (no cluster heterogeneity)
    \item[\textbf{A2$'$}] Within-cluster only ($\nuil$ may correlate with $\vii$); strictly weaker than A2
    \item[\textbf{A3}] Trajectory pooling: $E[\theta_i \mid \underline{d}_i, \vii] = E[\theta_i \mid \vii]$ (within cluster, comparative-advantage means do not vary across trajectories)
  \end{description}

\end{frame}

% ---------------------------------------------------------------
% Slide 15b: Three estimators, three assumption regimes (table + summary prose)
% finding 2: no \footnotesize
\begin{frame}{Three estimators, three assumption regimes}

  \begin{center}
  \begin{tabular}{lll}
    \toprule
    Estimator                                & Needs        & Targets      \\
    \midrule
    CKT main ($\texttt{run\_grc}$, Suri-style) & A1, A2       & $\phicoef$   \\
    CKT robust-vv ($\texttt{run\_grc\_robust\_vv}$) & A1, A2$'$, A3 & $\phicoef$   \\
    VV worker-level (robust)                 & A1, A2$'$    & $\alpha_1$   \\
    \bottomrule
  \end{tabular}
  \end{center}

  \vspace{0.4em}
  \begin{itemize}
    \item When A2 holds: all three target the same LCA relationship
    \item When only A2$'$ holds: CKT main is biased; CKT robust-vv and VV are consistent
    \item When A3 fails: CKT robust-vv picks up bias $\mathrm{Cov}_s(\bar\beta(s), \mud - \mu_{\underline{d}_0})$; VV stays consistent
  \end{itemize}

\end{frame}

% ---------------------------------------------------------------
% Slide 16: "Our implementation..."
\begin{frame}{Our implementation moves cluster structure to the instruments}

  \texttt{run\_grc\_robust\_vv}: keep CKT's parameters, swap the instruments
  \begin{itemize}
    \item Regress $\mathbbm{1}\{\underline{d}_i = s\}\,\D_{it}$ on $\mathbbm{1}\{\vii\}$ \emph{among trajectory-$s$ workers}
    \item Use residuals as instruments for $\phicoef$ identification
  \end{itemize}

  \vspace{0.4em}
  Same parameter set as the simple GRC ($\mud$, $\Delta_{\underline{d}_0}$, $\kappa$, $\gamma$); no $\beta(v)$ parameters

\end{frame}

% ---------------------------------------------------------------
% Slide 17: finding 6 --- replace "Pros:"/"Cons:" bold labels with
% two-column layout
\begin{frame}{The move buys a weaker assumption at tractable cost}

  \begin{columns}[t]
    \column{0.48\textwidth}
    Pros
    \begin{itemize}
      \item A2$'$ matches what the migration literature documents
      \item the CHN hukou split acts as a principled special case ($\vii = $ hukou) rather than an ad hoc fix
      \item stayer ATE estimates defensible against Verdier's concern in Suri's data
    \end{itemize}

    \column{0.48\textwidth}
    Cons
    \begin{itemize}
      \item $\Delta_{\dN}$ averages over provinces with switcher support, not the full population
      \item thinner identification of $\phicoef$: between-province switcher variation no longer contributes; only the within-province component does
      \item cluster-robust standard errors at the $\vii$ level
    \end{itemize}
  \end{columns}

\end{frame}

% ---------------------------------------------------------------
\section{Where the worker-level estimator can fail}

% ---------------------------------------------------------------
% Finding 3: split slide 18 into two slides
% Slide 18a: VV's setup
% finding 2: remove \footnotesize
\begin{frame}{VV uses two geographic levels}

  \citet{verdierAverageTreatmentEffects2020a}'s Suri-Kenya application uses two distinct geographic units:
  \begin{itemize}
    \item Province: enters as province$\times$year FE in Step 1 covariates (\S4.1)
    \item Village: cluster index $\vii$ for robust extrapolation (\S4.3); $\sim 94$ villages, $\sim 12$ farmers each
  \end{itemize}

  \vspace{0.6em}
  The two roles are distinct: province absorbs region-level trends in the outcome equation; village defines the within-cluster identification of $\phicoef$ in Step 2.

\end{frame}

% ---------------------------------------------------------------
% Slide 18b: our port and the fix
% finding 3: fold alertblock into plain prose
% finding 2: no \footnotesize
\begin{frame}{Our port clusters at the wrong geographic level}

  Three sub-province levels exist in our panels:
  \begin{itemize}
    \item CHN: \texttt{cid} (community), $\sim 4{,}066$ distinct, median $\sim 2$ workers
    \item IDN: \texttt{keca} (kecamatan / subdistrict), $\sim 1{,}669$ distinct, median $\sim 4$
    \item TZA: \texttt{ward}, $\sim 135$ distinct, median $\sim 97$
  \end{itemize}

  \vspace{0.6em}
  Our vanilla VV port currently clusters at first-wave province (\texttt{vfirst}, $\sim 25$ provinces).
  That choice came from a misreading: we saw ``province'' in VV's Step 1 and missed that his $\vii$ in Step 2 is the village.
  The fix is to rerun at the village-equivalent level (\texttt{cid} / \texttt{keca} / \texttt{ward}), keeping province as a Step 1 covariate per VV \S4.1.

\end{frame}

% ---------------------------------------------------------------
% Slide 19: finding 6 --- drop \textbf{Fragility cases:}; finding 7 --- drop
% trailing semicolons/periods from three bullets
\begin{frame}{Cluster-FE IV is fragile when within-cluster movers are scarce}

  Identification needs (i) within-cluster movers, (ii) trajectory variation among them, (iii) cluster size large enough for FE-IV to be well-conditioned.

  \vspace{0.4em}
  Fragility cases:
  \begin{itemize}
    \item zero movers in a cluster $\Rightarrow$ cluster-residualized period-treatment indicators are zero, no contribution to identification
    \item one or two movers $\Rightarrow$ within-cluster demeaning eats most of the variation
    \item with $T$ periods and $k$ movers, the within-cluster trajectory design has rank at most $\min(T, k-1)$ after partialling worker FEs; clusters with $k \le T$ contribute fewer free moments than nominal
  \end{itemize}

  \vspace{0.4em}
  This is a general property of cluster-FE IV with worker-level fixed effects.
  The binding question is not ``is $V$ too large?'' but ``how many clusters have too few movers?''

\end{frame}

% ---------------------------------------------------------------
% Slide 20: finding 11 --- reword title (drop "(pending)" hedge)
\begin{frame}{An observed shifter $w_i$ would provide a separate falsification test}

  Regress estimated $\thetai^{*}$ on $\Deltai$ and an observed shifter $w_i$:
  \[
    \thetai^{*} = \alpha_0 + \alpha_1 \Deltai + \alpha_2 w_i + e_i,
    \qquad H_0: \alpha_2 = 0
  \]
  Kenya: $w_i = $ distance to nearest seed seller; rejects if $\alpha_2 \neq 0$

  \vspace{0.5em}
  $w_i$ must be observed, exogenous to skill conditional on $\Deltai$, and correlated with $\nuil$

  \vspace{0.5em}
  Candidates in our data:
  \begin{itemize}
    \item distance / travel time to nearest urban center
    \item rainfall variability or pre-period climate shock (mergeable via GPS/admin code)
    \item distance to family already in urban area (closest analogue, but reflection problem)
  \end{itemize}

  \vspace{0.3em}
  None is perfect; data availability is the binding constraint

\end{frame}

% ---------------------------------------------------------------
% Finding 8: add takeaway slide before appendix
\begin{frame}{A2$'$ is the right assumption; clustering at the village level is the fix}

  \begin{itemize}
    \item A2$'$ is strictly weaker than A2 and matches what the migration literature documents---region-level cost heterogeneity is real, not a minor nuisance
    \item we currently cluster at the province level, which is VV's Step 1 covariate level, not his Step 2 cluster; the fix is to switch to the village-equivalent (\texttt{cid} / \texttt{keca} / \texttt{ward})
    \item once we rerun at the village-equivalent level, CKT-robust-vv identifies $\phicoef$ from within-community switcher variation and the extrapolation to non-migrants is defensible under A2$'$
  \end{itemize}

\end{frame}

% ---------------------------------------------------------------
\section{Appendix: parameter mapping}

% ---------------------------------------------------------------
\begin{frame}{Object-by-object map: worker and trajectory objects}

  \begin{center}
  \begin{tabular}{p{4.0cm}p{4.5cm}p{4.5cm}}
    \toprule
    Object & VV & CKT \\
    \midrule
    \multicolumn{3}{l}{\emph{Worker-level objects} (parameterized only in VV)} \\
    Worker rural mean                  & $\hat a_i$                            & not parameterized                                  \\
    Worker urban mean                  & $\hat a_i + \hat b_i$                  & not parameterized                                  \\
    Worker return                      & $\hat b_i$                            & not parameterized                                  \\
    \addlinespace
    \multicolumn{3}{l}{\emph{Trajectory-level objects} (CKT's natural level)} \\
    Trajectory rural mean              & $\bar{\hat a}_g$ (post-Step-1 avg in bin $g$) & $\mud$ (joint GMM parameter) \\
    Trajectory return                  & $\bar{\hat b}_g$                       & $\Deltad$ (joint GMM parameter) \\
    \bottomrule
  \end{tabular}
  \end{center}

\end{frame}

% ---------------------------------------------------------------
\begin{frame}{Object-by-object map: LCA structure}

  \begin{center}
  \begin{tabular}{p{4.0cm}p{4.5cm}p{4.5cm}}
    \toprule
    Object & VV & CKT \\
    \midrule
    \multicolumn{3}{l}{\emph{LCA structure}} \\
    LCA equation                       & $a_i = \alpha_0 + \alpha_1 b_i + e_i$  & $\Deltai = \beta + \phicoef\thetai + \xi_i$ \\
    Slope parameter                    & $\alpha_1$ (slope of $a$ on $b$)       & $\phicoef$ (slope of $\Delta$ on $\theta$) \\
    Right-hand-side variable           & rural baseline $a_i$                  & comparative advantage $\thetai$ \\
    \bottomrule
  \end{tabular}
  \end{center}

\end{frame}

% ---------------------------------------------------------------
% Finding 4: split into two appendix slides at \footnotesize
% Column widths tuned to Beamer 169 text width (~12.8cm usable):
% p{2.8cm}+p{3.5cm}+p{2.5cm}+p{2.8cm} = 11.6cm + tabcolsep * 6 ~ 12.6cm
\begin{frame}{Object-by-object map: geographic structure}

  \begin{center}
  \begin{tabular}{p{3.2cm}p{3.5cm}p{2.5cm}p{3.1cm}}
    \toprule
    Object & VV & CKT-main & CKT-vv \\
    \midrule
    Province (covariate use)
      & province$\times$year FE in Step 1
      & not used as a separate level
      & not used as a separate level \\
    Cluster index $\vii$
      & village (${\sim}94$, ${\sim}12$ workers each)
      & none
      & first-wave province (${\sim}25$); to be switched to \texttt{cid}/\texttt{keca}/\texttt{ward} \\
    \bottomrule
  \end{tabular}
  \end{center}

  \vspace{0.4em}
  {\footnotesize CKT-main: \texttt{run\_grc}. \quad CKT-vv: \texttt{run\_grc\_robust\_vv}.}

\end{frame}

% ---------------------------------------------------------------
\begin{frame}[allowframebreaks]{Object-by-object map: estimator details}

  \begin{center}
  \begin{tabular}{p{2.5cm}p{3.2cm}p{3.0cm}p{3.4cm}}
    \toprule
    Object & VV & CKT-main & CKT-vv \\
    \midrule
    Step 1 estimator
      & OLS w/ worker$\times$regime FE
      & n/a (joint GMM)
      & n/a (joint GMM) \\
    Step 2 estimator
      & IV: $\hat a$ on $\hat b$, cluster FE
      & joint GMM, switcher$\times$choice instruments
      & joint GMM, cluster-residualized switcher$\times$choice instruments \\
    Aggregation level
      & worker (movers only)
      & trajectory
      & trajectory \\
    Inference
      & cluster bootstrap or analytical, at $\vii$
      & default GMM SE
      & cluster-robust at $\vii$ \\
    \bottomrule
  \end{tabular}
  \end{center}

\end{frame}

% ---------------------------------------------------------------
% Finding 6: drop \textbf{Plain English.}, \textbf{Algebra.},
% \textbf{Two ways the identity breaks.} --- let prose carry the structure
% finding 2: remove \footnotesize
\begin{frame}{When the slopes coincide, and when they don't}

  Suppose every worker's return fell exactly on the LCA line (no noise) and there were one underlying skill so that rural baseline $a_i$ and comparative advantage $\thetai$ were the same object up to a constant.
  Then $\Deltai$ would be a perfect linear function of $a_i$ and the inverse would also be linear with reciprocal slope.

  \vspace{0.4em}
  In algebra: $\Deltai = \beta + \phicoef a_i$ inverts to $a_i = -\beta/\phicoef + (1/\phicoef)\Deltai$, so $\alpha_1 = 1/\phicoef$.

  \vspace{0.4em}
  Two ways this identity breaks:
  \begin{itemize}
    \item First, LCA noise: OLS slope of $a$ on $b$ is $\mathrm{Corr}(a,b)\,\sigma(a)/\sigma(b)$;
          slope of $b$ on $a$ is $\mathrm{Corr}(a,b)\,\sigma(b)/\sigma(a)$.
          Product equals $\mathrm{Corr}(a_i,b_i)^2 \le 1$.
          Reciprocal only when $|\mathrm{Corr}(a_i,b_i)| = 1$.
    \item Second, two-skill structure: $\thetai = b_R(\theta_i^U - \theta_i^R)$ uses both skills; $a_i$ uses $\theta_i^R$ alone.
          Slope of $\Delta$ on $\theta$ $\ne$ slope of $\Delta$ on $a$.
  \end{itemize}


\end{frame}

% ---------------------------------------------------------------
\begin{frame}[allowframebreaks]{References}
  \footnotesize
  \bibliographystyle{plainnat}
  \bibliography{../CKT.bib}
\end{frame}

% ---------------------------------------------------------------
\end{document}
